\documentclass[12pt, oneside]{book}

% --- PACKAGES ---
\usepackage{amsmath}       % Math equations
\usepackage{xcolor}        % Colored text
\usepackage{tcolorbox}     % Boxed text/equations
\usepackage{tikz}          % Drawing and image overlays
\usepackage{graphicx}      % Images
\usepackage[colorlinks=true, linkcolor=blue]{hyperref} % Interactive PDF links
\usepackage{parskip}
\usepackage[T1]{fontenc}
\usepackage{amssymb}
\usepackage{amsfonts}
\usepackage{bbm}

\begin{document}

\pagenumbering{roman}
\begin{titlepage}
    \centering
    
    \vspace*{3cm} 
    
    {\Huge \textbf{FEM Basics}\par}
    
    \vspace*{3cm} 
    \vfill
    
    {\small Oliver Jackson\par}
    
    \vspace*{0.5cm} 

    {\small \today\par}
\end{titlepage}


\tableofcontents

\cleardoublepage
\pagenumbering{arabic}

\chapter{Introduction}

This document details the FEM Galerkin method. 

\chapter{FEM basics}

Homework \#1: use the ID FEM to solve the following equation:

$$
-\frac{d}{d x}\left(e^{x} \frac{d u(x)}{d x}\right)=-e^{x}[\cos (x)-2 \sin (x)-x \cos (x)-x \sin (x)]
$$

For $0 \leq x \leq 1$ with $u(0)=0, u(1)=\cos (1)$. In standard FDM we apply the derivatives and then plug in our approximation for first and second order derivatives. In FEM it is advantageous to keep the lower order terms and convert into weak/variational form. It is important to keep a general form so it is easy to translate into code. Let

$$
c(x)=e^{x} \text { and } f(x)=-e^{x}[\cos (x)-2 \sin (x)-x \cos (x)-x \sin (x)]
$$

Then we are solving for

$$
-\frac{d}{d x}\left(c(x) \frac{d u}{d x}\right)=f(x) \quad \text { with solution } u
$$

The space $c^{k}(\bar{\Omega})$ is defined to be the set of all real-valved functions a defined on $\bar{\Omega}$ (which is both $\Omega$ and $\partial \Omega$ ) with the property that $u$ and its derivatives up to the $k^{\text {th }}$ order are all continuous on $\bar{\Omega}$. A solution to $u$ is sought in the subspace

$$
C_{D}^{2}(\bar{\Omega})=\left\{u \in C^{2}(\bar{\Omega}): u=0 \text { on } \partial \Omega\right\} .
$$

Introduce the test function $v$ and the space of test functions to be $C_{D}^{2}(\bar{\Omega})$.

$$
-\nabla(c \cdot \nabla u)=f
$$

$$
\int-\nabla(c \cdot \nabla u) \cdot v=\int f \cdot v
$$

This equation holds for some $u \in C_{D}^{2}(\bar{\Omega})$ denoted $u_{h}$ and for all $v \in C_{D}^{2}(\bar{\Omega})$ if and only if $u$ satisfies the original differential equation.

The next step is to apply green's identity to the left hand side. A quick review of Green's identity. The weak form of a BUP is derived using greens identity, which is the multi-dimensional form of integration by parts.

$$
\nabla \cdot(v \nabla u)=\nabla v \cdot \nabla u+v \Delta u
$$

Then to obtain Greens identity integrate both sides over $\Omega$ and apply the divergence theorem

$$
\begin{aligned}
& \Rightarrow \int_{\Omega} \nabla \cdot(v \nabla u)=\int_{\Omega} \nabla v \cdot \nabla u+\int_{\Omega} v \Delta u \\
& \Rightarrow \int_{\partial \Omega} v \nabla u \cdot n=\int_{\Omega} \nabla v \cdot \nabla u+\int_{\Omega} v \Delta u \\
& \Rightarrow-\int_{\Omega} v \Delta u=\int_{\Omega} \nabla v \cdot \nabla u-\int_{\partial \Omega} v \frac{\partial u}{\partial n}
\end{aligned}
$$

we know $V=0$ on $\partial \Omega$ so

$$
-\int_{\Omega} v \Delta u=\int_{\Omega} \nabla V \cdot \nabla u
$$

So we get

$$
\int_{\Omega}-\nabla(c \cdot \nabla u) \cdot v=\int_{\Omega} c \nabla u \cdot \nabla v
$$

So the weak form of our elliptic problem is given by Find $u \in C_{D}^{2}(\bar{\Omega})$ such that $\int_{\Omega} C \nabla u \cdot \nabla_{v}=\int_{\Omega}$ fv for all $v \in C_{D}^{2}(\bar{\Omega})$ using notation $(a, b)=\int_{\Omega} a b d x$ and bilinear operator $a$,

$$
a(u, v)=(f, v) \quad \forall v \in C_{D}^{2}(\bar{\Omega})
$$

Where $a(u, v) \equiv \int_{\Omega} c \nabla u \cdot \nabla v$. Define nodes on the interval $x_{0}<x_{1}<\cdots<x_{N}<x_{N+1}$. The space between the nodes $N$ are defined $h=x_{n}-x_{n-1}$ where $n=0, \ldots, N+1$

The space $C_{D}^{2}(\bar{\Omega})$ has linear or quadratic functions on each interval and we may choose the values $M_{n}=V\left(X_{n}\right)$ at the node points $n=0, \ldots, N+1$. We can introduce basis functions such that overlapping and neighboring hodes return a value and disconnected nudes return 0 .

$$
\phi_{j}\left(x_{i}\right)= \begin{cases}1 & \text { if } \quad i=j \\ 0 & \text { if } \quad i \neq j \quad i, j=1, \ldots, N\end{cases}
$$

which can be linear or quadratic. A function $V \in C_{D}^{2}(\bar{\Omega})$ then has the representation

$$
V(x)=\sum_{m=1}^{M} \eta_{n} \phi_{n}(x) \quad x \in[0,1]
$$

Then the FEM can now be formulated as: Find $u_{h} \in C_{D}^{2}(\bar{\Omega})$ such that

$$
a\left(u_{h}, v\right)=(f, v) \quad \forall v \in C_{D}^{2}(\bar{\Omega})
$$

If $u_{n}$ satisfies the above then

$$
a\left(u_{h}, \phi_{n}\right)=\left(f, \phi_{n}\right) \quad n=1, \ldots, N \not \otimes
$$

Then by taking linear combinations we see un satisfies \& since

$$
u_{h}(x)=\sum_{n=1}^{N} \xi_{m} \phi_{m}(x), \quad \xi_{m}=u_{h}\left(x_{m}\right)
$$

rewrite * like so

$$
\sum_{n=1}^{N} \xi_{n} a\left(\phi_{n}, \phi_{m}\right)=\left(f, \phi_{n}\right) m=1, \ldots, M
$$

which is a linear system of equations with M equations and $N$ unknowus $\xi_{1}, \ldots, \xi_{N}$ which we are solving for. In matrix form the linear system of equations can be written as

$$
A \xi=b
$$

Where $A=\left(A_{m n}\right)$ is the $M \times N$ matrix with elements $A_{m n}=a\left(\phi_{n}, \phi_{m}\right)$ called stiffness matrix. And where $\xi=\left(\xi_{1}, \ldots, \xi_{N}\right)$ and $b=\left(b_{1}, \ldots, b_{N}\right)$ with $b_{i}=\left(f, \phi_{i}\right)$ are both column vectors. $b$ is the load vector. Row $m$ is test and column $n$ is trail.

Now we want to build a general algorithm to construct the stiffness matrix which decouples the geometric shape of the mesh from the functional degrees of freedom (the nodes).

The algorithm will be general in the sense that any number of nodes $N$ can be chosen and the geometry, dimension, and basis type are separated through if /else statements in the code. Forther calculus and Jacobian matrix is required if a general function for the construction of mesh element and nodes. Instead the general framework advised by Prof He is to add additional matrix formulations for linear VS quadratic basis functions into if/else statement.

The $P$ matrix is an information matrix consisting of the coordinates of all mesh nodes. It is a $D \times N$ matrix where $D$ is dimension. In $1 D$ and $2 D$

$$
\begin{aligned}
P_{1 D} & =\left(x_{1}, \ldots, x_{N}\right) \\
P_{2 D} & =\binom{x_{1}, \ldots, x_{N}}{y_{1}, \ldots, y_{N}}
\end{aligned}
$$

The $T$ matrix is an information matrix consisting of the global node indices of the mesh elements. The elements are the integral bounds. So the $T$ matrix holds the indices of the elements $E_{k}$ where $k$ denotes the number of elements $k=1, \ldots, N-1$. Define $B$ to be the number of basis functions required to construct one element. Then $T$ is a $B X K$ matri $x$. For the line ar basis function case

$$
T=\left(\begin{array}{lllll}
1 & 2 & \cdots & N-2 & N-1 \\
2 & 3 & \cdots & N-1 & N
\end{array}\right)
$$

when we have linear basis functions define
contains $\Omega$ points $P=P_{b}$ where $b$ denotes basis node value contains $\Omega$ elements $T=T_{b}$ where $b$ denotes basis element value

So when our domain is defined with quadratic basis functions then the number of integrals we take increases and this also the number of $a_{n m}$ values increases but NOT the number of elements in our domain. Thus $P$ and $T$ contain information about the geometry of the system which is represented in the element stiffness matrix and $\quad T_{b}$ contains algebraic information used to map the element basis stiffness matrix to the global stiffness matrix. Pb contains algebraic information about basis function construction. For quadratic
basis $P \neq P_{b}$ and $T \neq T_{b}$.
Now consides the interval of a specific element $E=\left[x_{n}, x_{n+1}\right]$. Then because our basis functions are defined for $i \neq j \quad \phi_{i}\left(x_{j}\right)=0$ and $\phi_{i}\left(x_{i}\right)=1$, the only basis functions we consides over this element $E$ are $\phi_{n}, \phi_{n+1}$. Note the following important property

$$
\left(\phi_{j}, \phi_{j+1}\right)=\left(\phi_{j+1}, \phi_{j}\right)
$$

Then we define local basis function $\psi$ as $\phi$ over the interval E.
Recall the linear basis functions:

$$
\phi_{n}(x)=\left\{\begin{array}{cl}
\frac{x-x_{n-1}}{h} & \text { for } x \in\left[x_{n-1}, x_{n}\right] \\
\frac{x_{n+1}-x}{h} & \text { for } x \in\left[x_{n}, x_{n+1}\right] \\
0 & \text { else }
\end{array}\right.
$$

Then define the local basis functions for an element $E=\left[x_{n}, x_{n+1}\right]$. On this element we have contributions from the positive edge of $\phi_{n+1}$ and negative from $\phi_{n}$ call these slopes $\psi_{n}, \psi_{n+1}$ given by

$$
\begin{aligned}
& \psi_{1}=\frac{x_{n+1}-x}{h} \text { for } x \in\left[x_{n}, x_{n+1}\right] \\
& \psi_{2}=\frac{x-x_{n}}{h} \text { for } x \in\left[x_{n}, x_{n+1}\right]
\end{aligned}
$$

So in one element $E_{k}$ we have two local basis functions denoted $N_{b}=2$ for linear global basis functions and the total number of global basis functions is $N$. We need two points $\left[x_{n}, x_{n+1}\right]$ to compute $\psi_{1}, \psi_{2}$ that are extracted
directly from the $P_{b}$ matrix using the indices from the Tb matrix:

$$
\begin{aligned}
& x_{n}^{(k)}=P_{b}\left(1, T_{b}(1, k)\right) \\
& x_{n+1}^{(k)}=P_{b}\left(1, T_{b}(2, k)\right)
\end{aligned}
$$

There are four non-zero local integrals on $E_{k}=\left[x_{n}, x_{n+1}\right]$

$$
\begin{aligned}
& a_{\alpha \beta}^{k}=a^{k}\left(\psi_{\alpha}, \psi_{\beta}\right) \text { for } \alpha, \beta \in\{1,2\} \\
& a_{\alpha \beta}^{k}=\int_{x_{n}^{k}}^{x_{n+1}^{k}} c\left(\psi_{\alpha}^{\prime} \psi_{\beta}^{\prime}\right) d x \text { for } \alpha, \beta \in\{1,2\}
\end{aligned}
$$

This means the local integral should be assembled to the global matrix entry $a_{m n}$. An element is given by $E_{k}=\left[x_{n}^{k}, x_{n+1}^{k}\right]$ for $k=1, \ldots, N-1$. We compute an $N_{b} \times N_{b}$ element stiffness matrix $a^{k}$ with indices $\alpha, \beta \in\{1,2\}$ for the linear case. Because $\beta$ is the test function index it indicates the rows of $a^{k}$ and $\alpha$ the columns.

$$
a_{\alpha \beta}^{k}=\left[\begin{array}{cc}
a_{11}^{k} & a_{12}^{k} \\
a_{21}^{k} & a_{22}^{k}
\end{array}\right]=\left[\begin{array}{ll}
\int\left(\psi_{1}^{\prime} \psi_{1}^{\prime}\right. & \int c \psi_{1}^{\prime} \psi_{2}^{\prime} \\
\int c \psi_{2}^{\prime} \psi_{1}^{\prime} & \int c \psi_{2}^{\prime} \psi_{2}^{\prime}
\end{array}\right]
$$

To map the local element stiffness matrix $a_{\alpha \beta}^{k}$ into the global $M \times N$ matrix $A$ we use the topology matrix $T_{b}$ to route the indices.

$$
\begin{aligned}
& \text { Global cow m }=T_{b}(\beta, k) \\
& \text { Global column } n=T_{b}(\alpha, k)
\end{aligned}
$$

I write the global entry of $A$ as $A_{m n}$ and the local element stiffness matrix as $a_{\beta \alpha}^{k}$. I include $\delta_{m, T_{b}(\beta, k)}$ as a way to search for the global row index and $\delta_{n}, T_{b}(\alpha, k)$ as a way to search for the global column index.

$$
A_{m n}=\sum_{n=1}^{N-1} \sum_{\beta=1}^{N_{b}} \sum_{\alpha=1}^{N_{b}} a_{\alpha \beta}^{k} \cdot \delta_{m, T_{b}(\beta, k)} \cdot \delta_{n, T_{b}(\alpha, k)}
$$

To program this we can take the steps

NOTATION KEY

- $N=$ number of internal nudes in domain $N$ which corresponds to the \# of cukhowns
- $k=G$ lobal element index $(k=1, \ldots, N-1) k$
- $m=$ Global test basis index $(1, \ldots, M)$ maps rows of $A M$
- $n=$ Global trail basis index $(1, \ldots, N)$ maps to columns of $A n$
- $\beta=$ local test basis index $\beta \in\{1,2\}$ beta
- $\alpha=$ local trail basis index $\alpha \in\{1,2\}$ alpha
- $\xi_{n}=$ global trial coefficients $x_{i}$
- $\eta_{m}=$ global test coefficients eta
- $a_{\beta \alpha}^{k}=$ local stiffness matrix local-a
- domain $=X=x_{0}, \ldots, x_{N+1} \quad x$
- $T, T_{b}=$ topology information matrix $T$-matrix, $T$ - $b$-matrix
- $P, P_{b}=$ index information matrix $P$-matrix, $P$ - $b$-matrix

\end{document}
